\chapter{The Input System}
\label{ch:input-system}

CNA's \cnans{Microsoft::Xna::Framework::Input} namespace covers keyboard, mouse, gamepad,
and touch --- and, layered underneath it, a set of CNA-only extensions with no XNA
counterpart at all. Its own audit (\texttt{plan\_input.md}, 505 tasks across 13 phases) is
fully closed: 490 done, 15 formally blocked (all real-hardware-only checks that cannot be
automated at all), zero open.

\section{Architecture: event-driven, not poll-driven}

The most important structural fact about this namespace, and the one place CNA's input
model genuinely diverges from FNA's own: FNA re-queries SDL fresh every time
\texttt{GetState()} is called; CNA instead accumulates input state as SDL events are pumped,
and \texttt{GetState()} is a pure read of that accumulated snapshot. The pipeline runs in one
direction: \texttt{Game::PollEvents()} calls \texttt{SDL\_PollEvent} once per frame, funneling
every event through a single dispatcher --- \texttt{SdlInputBridge::ProcessEvent} --- which
updates \texttt{InputManager}'s internal state, \texttt{TextInputEXT}'s internal state, and
\cnaclass{TouchPanel}'s gesture detector as appropriate; \texttt{Keyboard} / \texttt{Mouse} /
\texttt{GamePad} / \texttt{TouchPanel::GetState()} then just read that snapshot back.
\texttt{FrameworkDispatcher::Update()} (this book's game-loop chapter) drives
\cnaclass{TouchPanel}'s own per-frame gesture update. For ordinary once-per-frame calling
code this is behaviorally equivalent to FNA's approach; the one place it is not is calling
\texttt{GetState()} twice within the same frame, which returns two identical snapshots in
CNA where FNA might observe a change in between. This entire pipeline is deliberately
single-threaded and unsynchronized --- safe only because both halves (the event pump and the
\texttt{Update} / \texttt{Draw} reads) run on the same game-loop thread, matching both SDL's own
requirement that events be pumped on the video thread and XNA/FNA's own single-threaded
assumption.

Three tags recur through this namespace's own headers and documentation, and are worth
knowing before reading further: \textbf{STRICT} (must match FNA's real surface exactly),
\textbf{EXT} (an FNA-compatible extension, using the \texttt{...EXT} suffix convention from
this book's game-loop chapter), and \textbf{NOXNA} (a CNA-only addition with no FNA counterpart
at all). A dedicated compile-time test suite (\texttt{PublicApiInputSignatureFreezeTests.cpp})
pins the exact strict-XNA member set this namespace promises never to silently change.

\section{Keyboard}

\cnaclass{Keyboard} is static-only, offering \texttt{GetState()} and a
per-\cnaclass{PlayerIndex} overload. \cnaclass{KeyboardState} deliberately deviates from
FNA's own storage layout --- backed by a \texttt{std::unordered\_set<Keys>} rather than
FNA's native bitfield --- on the reasoning that the observable \emph{contract}
(\texttt{IsKeyDown}, \texttt{GetPressedKeys()} returning ascending-sorted keys, a specific
\texttt{GetHashCode()} formula reconstructing FNA's own eight-32-bit-word XOR layout) is what
must match, not the internal representation producing it. \cnaclass{Keys} itself has over
150 values, independently confirmed byte-for-byte identical to FNA's numeric values,
including several distinctly odd hexadecimal outliers (\texttt{Pause = 0x13},
\texttt{Kana = 0x15}, \texttt{ChatPadGreen = 0xCA}) preserved exactly rather than renumbered
for tidiness. A real, load-bearing platform limit: forty \cnaclass{Keys} values have no SDL
scancode source at all --- IME keys, ChatPad keys, and most browser/media keys other than
volume up/down --- and simply cannot round-trip through the \texttt{EXT} scancode-lookup
helper, matching an identical omission in FNA's own key/scancode maps rather than being a
CNA-specific gap.

\subsection{Full method reference: KeyboardState and Keyboard's EXT surface}

\cnaclass{KeyboardState}'s own public surface, read directly from its header, is smaller than
its internal \texttt{std::unordered\_set<Keys>} storage might suggest:

\begin{description}
\item[\texttt{KeyboardState(std::initializer\_list<Keys> keys)}] the ordinary way application
  and test code constructs one directly --- \texttt{KeyboardState\{Keys::W, Keys::LeftShift\}}
  for a manually-composed state.
\item[\texttt{getItem(Keys key)} / \texttt{operator[](Keys key)}] both return a
  \cnaclass{KeyState} (\texttt{Down} or \texttt{Up}) for one key --- the indexer form XNA
  itself exposes as \texttt{keyboardState[Keys.W]}.
\item[\texttt{IsKeyDown(key)} / \texttt{IsKeyUp(key)}] the ordinary boolean queries --- what
  almost all application code actually calls.
\item[\texttt{GetPressedKeys()}] returns every currently-down key as an ascending-sorted
  \texttt{std::vector<Keys>} --- the ascending-sort guarantee is explicit and load-bearing:
  it is part of the observable contract this chapter's own intro named as what CNA
  deliberately preserves even though the backing storage (a hash set) has no inherent order
  of its own.
\end{description}

Beyond \cnaclass{KeyboardState} itself, \cnaclass{Keyboard} (the static entry point) exposes a
cluster of \texttt{EXT} methods with no FNA equivalent, all revolving around one real gap
XNA's own \cnaclass{Keys} enum has: it is fundamentally a \emph{layout-dependent} enum (its
values correspond to US-QWERTY key positions and their typical labels), while SDL3 itself
tracks both a layout-independent \emph{scancode} (physical key position) and a
layout-\emph{dependent} keycode. \texttt{GetKeyFromScancodeEXT(scancode)} and its inverse
direction let a caller translate between the two --- the mechanism a non-US keyboard layout
needs to implement ``WASD, but by physical position rather than by label'' correctly.
\texttt{GetModStateEXT()} returns the live modifier-key bitmask (shift/ctrl/alt/etc.)
\cnaclass{Keys} alone cannot represent as a single down/up state per key.
\texttt{GetScancodeNameEXT} / \texttt{GetKeyNameEXT} (and their name-to-key inverses) return a
human-readable, current-layout-aware label for a key or scancode --- what a ``press any key to
continue'' or key-rebinding UI should actually display, since \texttt{Keys::W} should read as
``Z'' on an AZERTY layout, not literally ``W''.

\begin{lstlisting}
// Worked example: physical-position movement keys that stay correct on a
// non-QWERTY layout -- the real reason GetKeyFromScancodeEXT exists.
// WASD's physical positions, expressed as scancodes rather than as the
// QWERTY-specific Keys::W/A/S/D values directly.
const Keys forwardKey = Keyboard::GetKeyFromScancodeEXT(Keys::W);
const Keys leftKey    = Keyboard::GetKeyFromScancodeEXT(Keys::A);
const Keys backKey     = Keyboard::GetKeyFromScancodeEXT(Keys::S);
const Keys rightKey   = Keyboard::GetKeyFromScancodeEXT(Keys::D);

// On an AZERTY keyboard, forwardKey now resolves to the physical key an
// AZERTY user actually presses for "forward" (labelled Z, not W) --
// GetKeyFromScancodeEXT(Keys::W) treats Keys::W as identifying the
// physical scancode a US layout would call W, then maps that scancode
// back through the *current* layout to find what key is physically there.
const KeyboardState keyboard = Keyboard::GetState();
if (keyboard.IsKeyDown(forwardKey))
{
    // move forward
}
\end{lstlisting}

\section{Mouse}

\cnaclass{Mouse} is static-only: \texttt{GetState()} and \texttt{SetPosition(x, y)}, the
latter carefully converting logical coordinates to real window coordinates before calling
SDL's warp function --- and doing so differently per backend (offset-aware on SDL\_Renderer,
uniform-scale-aware on EasyGL, pass-through on Vulkan/BGFX), a small, concrete illustration of
how even an ostensibly backend-agnostic input call has to account for
this book's per-backend presentation differences (its backend-architecture chapter).
\cnaclass{MouseState} carries the expected X/Y, five buttons, and a cumulative,
120-units-per-notch \texttt{ScrollWheelValue}; a \texttt{NOXNA} horizontal-scroll addition is
deliberately excluded from \texttt{Equals} / \texttt{GetHashCode} / \texttt{ToString} specifically
so those stay byte-identical to FNA's own output despite the extra field existing. A wholly
\texttt{NOXNA} class, \cnaclass{MouseCursor}, wraps \texttt{SDL\_Cursor*} with eleven lazily
created, process-lifetime stock system cursors --- a case where CNA's design is deliberately
\emph{more} defensive than the MonoGame implementation it is closest to: MonoGame's own
\texttt{PlatformDispose()} frees a shared stock cursor handle unconditionally, which can
corrupt the shared handle if any other code still holds it; CNA makes \texttt{Dispose()} a
safe no-op for the shared stock instances specifically to avoid that failure mode.

\subsection{Full method reference: MouseState and Mouse's EXT surface}

\cnaclass{MouseState} is a plain data carrier --- \texttt{X} / \texttt{Y}, five
\cnaclass{ButtonState} properties (\texttt{LeftButton} / \texttt{RightButton}/%
\texttt{MiddleButton} / \texttt{XButton1} / \texttt{XButton2}), the cumulative
\texttt{ScrollWheelValue} already named above, and (\texttt{NOXNA}) a parallel
\texttt{HorizontalScrollWheelValueEXT} for horizontal scroll, deliberately excluded from
\texttt{Equals} / \texttt{GetHashCode} / \texttt{ToString} as already noted. \cnaclass{Mouse}
itself, beyond \texttt{GetState()} and \texttt{SetPosition}, exposes a real, working relative
(pointer-lock-style) mouse mode and a small cluster of window/capture-related \texttt{EXT}
methods:

\begin{description}
\item[\texttt{IsRelativeMouseModeEXT}] (get/set)
  toggles SDL's own relative mouse mode, where motion is reported as unbounded deltas rather
  than an absolute, screen-clamped position --- the standard technique for an FPS-style camera
  look. Both accessors are safe to call with no window at all (no published window handle and
  no focused window): the getter simply reports \texttt{false} rather than querying SDL with a
  null window (whose behavior SDL itself leaves undefined), and the setter is a silent no-op
  for the same reason. Reading \texttt{InputManager::GetMouseState()} directly shows precisely
  how the substitution works: every \texttt{SDL\_EVENT\_MOUSE\_MOTION} event accumulates into
  a separate \texttt{RelativeDeltaX} / \texttt{RelativeDeltaY} pair regardless of mode, but
  \texttt{MouseState}'s public \texttt{X} / \texttt{Y} only read from that accumulator --- and
  reset it to zero immediately after --- while relative mode is active; otherwise \texttt{X}/%
  \texttt{Y} report the ordinary absolute position. This is why \texttt{MouseState.X} / \texttt{Y}
  genuinely mean two different things depending on \texttt{IsRelativeMouseModeEXT}, not just a
  documentation convention.
\item[\texttt{SetPosition(x, y)}] is itself a documented no-op while relative mouse mode is
  active --- read directly from \texttt{Mouse.cpp}: the method's very first check returns
  immediately if \texttt{getIsRelativeMouseModeEXTProperty()} is true, before ever calling
  SDL's warp function, since an absolute position warp is meaningless once the mouse is
  reporting relative deltas instead.
\item[\texttt{SetCaptureEXT(enabled)}] engages OS-level mouse capture (input keeps arriving
  even if the cursor leaves the window bounds) --- distinct from relative mode, and combinable
  with it.
\item[\texttt{GetGlobalPositionEXT} / \texttt{WarpGlobalEXT}] read or set the mouse position in
  \emph{desktop} (not window-relative) coordinates --- the form a multi-window or
  multi-monitor application needs, where an ordinary \texttt{GetState()} position is
  meaningless across window boundaries.
\item[\texttt{ClickedEXT}] a \texttt{NOXNA System::MulticastAction<int>} event (Chapter~34
  covers this \texttt{sharp-runtime} type in full) firing once per real click with the button
  index, for code that wants a discrete click event rather than polling button state every
  frame.
\end{description}

\begin{lstlisting}
// Worked example: an FPS-style mouse-look camera using relative mouse
// mode -- SetPosition would be silently ineffective here, which is
// exactly the point: relative mode reports motion as deltas instead.
game.setIsMouseVisibleProperty(false);
Mouse::setIsRelativeMouseModeEXTProperty(true);

// Each frame: read accumulated relative motion, not an absolute position.
const MouseState mouse = Mouse::GetState();
const float yaw   = mouse.X * mouseSensitivity;
const float pitch = mouse.Y * mouseSensitivity;
cameraRotation = Quaternion::CreateFromYawPitchRoll(yaw, pitch, 0.0f) * cameraRotation;
\end{lstlisting}

\section{GamePad}

The full XNA surface --- \texttt{GetCapabilities}, \texttt{GetState} with configurable
\cnaclass{GamePadDeadZone} handling, \texttt{SetVibration} --- is wired to real
\texttt{SDL\_Gamepad} hardware, not stubbed. Dead-zone constants match FNA exactly
(\texttt{LeftDeadZone = 7849/32768}, \texttt{RightDeadZone = 8689/32768},
\texttt{TriggerThreshold = 30/255}). Beyond the strict XNA surface, an extensive
\texttt{EXT} layer exposes real hardware capability FNA itself already extends beyond stock
XNA: GUID, light-bar and trigger-rumble control, gyroscope/accelerometer reads, and several
CNA-only additions with no FNA counterpart at all (player-index assignment, power info,
button-label queries, device name/path/serial/firmware, Steam handle, connection state,
touchpad finger tracking).

\subsection{Real bugs found and fixed}

Several concrete, historically real bugs are worth naming. \texttt{SDL\_INIT\_GAMEPAD} was
never actually initialized, meaning zero gamepad events were ever delivered until fixed.
\texttt{GetCapabilities()} itself used to \emph{cancel active rumble} as a side effect ---
it probed rumble support by calling \texttt{SDL\_RumbleGamepad(0, 0, 0)}, which is itself a
stop-vibration call; fixed by switching to a non-mutating capability property query instead.
A stick-normalization bug divided the negative half of an analog stick's range by 32768
rather than FNA's own 32767 for \emph{both} halves, meaning every non-endpoint negative
sample diverged slightly from FNA's expected value (for example, a raw value of $-16384$
producing $-0.5$ in CNA against FNA's $-0.50001$) --- fixed to match FNA's asymmetric divisor
exactly, a good small example of how ``obviously correct-looking'' code can still diverge
from a reference implementation's own specific numerical quirks.

\subsection{Full method reference: GamePadState and the three GamePadDeadZone modes}

\cnaclass{GamePadState} composes four sub-structures, each a plain value carrier read directly
from its header: \cnaclass{GamePadThumbSticks} (\texttt{Left} / \texttt{Right}, each a
\cnaclass{Vector2}), \cnaclass{GamePadTriggers} (\texttt{Left} / \texttt{Right}, each a scalar
\texttt{float}), \cnaclass{GamePadButtons} (eleven \cnaclass{ButtonState} properties --- every
face/shoulder/stick-click/back/start/big button --- plus a \texttt{ButtonStateFromFlag(Buttons)}
convenience for looking one up by the \cnaclass{Buttons} flag enum instead of by name), and
\cnaclass{GamePadDPad} (\texttt{Up} / \texttt{Down} / \texttt{Left} / \texttt{Right}, each a
\cnaclass{ButtonState}). \cnaclass{GamePadState} itself adds \texttt{IsConnected},
\texttt{PacketNumber} (increments whenever the underlying hardware state actually changes ---
the standard XNA technique for detecting ``did anything change since I last checked'' without
comparing every field by hand), and \texttt{IsButtonDown} / \texttt{IsButtonUp} convenience
wrappers over the same \cnaclass{Buttons} flag enum.

The dead-zone handling is worth reading precisely rather than assuming all three
\cnaclass{GamePadDeadZone} modes merely differ in threshold size --- they differ in
\emph{shape}. Reading \cnaclass{GamePadThumbSticks}'s constructor directly: \texttt{None}
applies no filtering at all. \texttt{IndependentAxes} --- the one-argument
\texttt{GamePad::GetState}'s own default choice --- excludes each axis's dead zone
independently, one call to \texttt{ExcludeAxisDeadZone} per axis: a stick pushed diagonally
near the dead-zone boundary can end up with one axis zeroed while the other still reports a
nonzero value, a small directional distortion right at the boundary. \texttt{Circular} instead
measures the stick's combined magnitude and excludes the dead zone radially, avoiding that
per-axis distortion at the cost of the extra trigonometry. Both non-\texttt{None} modes also
differ in \emph{clamp shape} after the dead zone is applied: \texttt{Circular} clamps the
result to a unit \emph{circle}; \texttt{None} and \texttt{IndependentAxes} both clamp to a unit
\emph{square} instead (independently per axis) --- so switching dead-zone modes can change a
stick's maximum diagonal magnitude, not merely its small-input behavior.

\begin{lstlisting}
// Worked example: dead-zone-aware movement plus rumble feedback, choosing
// Circular explicitly for the smoother diagonal response it gives over
// the IndependentAxes default -- most noticeable exactly at the boundary
// this book's own reading of GamePadThumbSticks's constructor identified.
const GamePadState pad = GamePad::GetState(PlayerIndex::One, GamePadDeadZone::Circular);
if (pad.getIsConnectedProperty())
{
    const Vector2 moveInput = pad.getThumbSticksProperty().getLeftProperty();
    playerPosition += moveInput * moveSpeed * deltaSeconds;

    if (pad.IsButtonDown(Buttons::A))
    {
        // Rumble briefly on jump -- SetVibration takes independent
        // left/right motor strengths in [0, 1], matching real Xbox 360
        // controllers' two distinct (low-frequency/high-frequency) motors.
        GamePad::SetVibration(PlayerIndex::One, 0.0f, 0.6f);
    }
}
\end{lstlisting}

Beyond the strict surface, a handful of the \texttt{EXT} methods named earlier are worth a
second, more concrete look: \texttt{GetGyroEXT} / \texttt{GetAccelerometerEXT} expose a modern
controller's own built-in motion sensors (distinct from a host device's own sensors, covered
later in this chapter) --- real only on hardware that actually reports them, and
\texttt{GamePadCapabilities::getHasGyroEXTProperty()} / \texttt{getHasAccelerometerEXTProperty()}
exist specifically so calling code can check before relying on either. \texttt{SetLightBarEXT}
and \texttt{SetTriggerVibrationEXT} (independent per-trigger rumble motors, distinct from the
two whole-controller motors \texttt{SetVibration} drives) are both real, modern-controller-only
capabilities with no XNA-era precedent at all, gated the same way behind their own
\texttt{GamePadCapabilities} flags.

\section{Touch and gestures}

\cnaclass{TouchPanel} is static-only (\texttt{MAX\_TOUCHES = 8}), providing
\texttt{GetCapabilities()}, \texttt{GetState()}, and \texttt{ReadGesture()}.
\cnaclass{TouchCollection} is immutable only by contract, not by construction --- its
\texttt{IsReadOnly} flag is advisory, matching an identical inconsistency already present in
FNA's own implementation. \cnaclass{GestureType} is an eleven-value bit-flag enum (Tap,
DoubleTap, Hold, Drag, Flick, Pinch, and their variants).

Several real bugs were found and fixed here too. \texttt{SDL\_EVENT\_FINGER\_CANCELED} was
simply unhandled, leaving a canceled touch stuck as if still active --- fixed to release it
the same way a genuine finger-up event would. \texttt{TouchPanel::GetCapabilities()} used to
consult only a sticky ``has this ever been touched'' flag rather than real device
enumeration, meaning a genuinely present but not-yet-touched touchscreen incorrectly reported
as disconnected on every non-Windows platform --- fixed to query the real SDL device list on
every call, keeping the sticky flag only as a fallback. And a frame-accuracy bug had touch
state advancement (promoting a touch from Pressed to Moved, retiring a Released touch) happen
\emph{inside} \texttt{GetState()} itself, meaning the number of times a frame happened to
call \texttt{GetState()} changed observable behavior --- fixed by splitting the operation
into a pure state read plus a separate, genuinely once-per-frame
\texttt{AdvanceTouchFrame()} step.

\subsection{Full method reference: TouchCollection, TouchLocation, and GestureSample}

\cnaclass{TouchCollection} (returned by \texttt{TouchPanel::GetState()}) is a fixed, per-frame
snapshot of every currently-tracked touch, read directly from its header:
\texttt{Count} / \texttt{IsConnected} (whether a touch device exists at all, distinct from
whether any finger is currently down), the by-index \texttt{operator[]}, and the same general
ordered-container surface Chapter~7's own \cnaclass{CurveKeyCollection} coverage already
described --- \texttt{Contains} / \texttt{IndexOf}, a \texttt{NOXNA} \texttt{empty()}, and
\texttt{NOXNA} iterator support for range-based \texttt{for}. Each \cnaclass{TouchLocation} inside it carries
an \texttt{Id} (stable across frames for the same physical finger, the correct key for
tracking one finger's motion over time), a \texttt{State} (\cnaclass{TouchLocationState}:
\texttt{Pressed}, \texttt{Moved}, \texttt{Released}, or \texttt{Invalid} for a stale/expired
entry), a \texttt{Position}, and (\texttt{NOXNA}) a \texttt{PressureEXT} reading on hardware
that reports one. \cnaclass{GestureSample} (returned by \texttt{TouchPanel::ReadGesture()})
carries a \texttt{GestureType}, a \texttt{Timestamp}, and up to two positions/deltas
(\texttt{Position} / \texttt{Position2}, \texttt{Delta} / \texttt{Delta2}) --- the second pair used
only by two-finger gestures like \texttt{Pinch}, left at their default for single-finger
gestures.

\begin{lstlisting}
// Worked example: per-finger drag tracking using TouchLocation's stable
// Id -- the field that makes tracking *one* finger's motion possible,
// since a touch's index within the TouchCollection is not stable
// frame-to-frame the way its Id is.
const TouchCollection touches = TouchPanel::GetState();
for (const TouchLocation& touch : touches)
{
    switch (touch.getStateProperty())
    {
    case TouchLocationState::Pressed:
        activeDrags[touch.getIdProperty()] = touch.getPositionProperty();
        break;
    case TouchLocationState::Moved:
        if (auto it = activeDrags.find(touch.getIdProperty()); it != activeDrags.end())
        {
            const Vector2 delta = touch.getPositionProperty() - it->second;
            PanCamera(delta);
            it->second = touch.getPositionProperty();
        }
        break;
    case TouchLocationState::Released:
    case TouchLocationState::Invalid:
        activeDrags.erase(touch.getIdProperty());
        break;
    }
}

// Gestures are read separately from raw touches, and only the types
// enabled via EnabledGestures are ever produced by ReadGesture().
TouchPanel::setEnabledGesturesProperty(GestureType::Pinch | GestureType::Flick);
while (TouchPanel::getIsGestureAvailableProperty())
{
    const GestureSample gesture = TouchPanel::ReadGesture();
    if (gesture.getGestureTypeProperty() == GestureType::Pinch)
    {
        // Position/Position2 are the two fingers' current positions;
        // Delta/Delta2 (unused by single-finger gestures) hold their
        // per-frame motion -- combine both to derive a zoom factor.
        ZoomCamera(gesture.getPositionProperty(), gesture.getPosition2Property());
    }
}
\end{lstlisting}

Two details worth calling out precisely, both readable directly from the enum/property
definitions rather than assumed: \texttt{EnabledGestures} is a genuine filter, not merely
advisory --- \texttt{ReadGesture()} only ever returns gesture types the caller has opted into,
so a UI that never enables \texttt{GestureType::Pinch} will never see one, even on hardware
that fully supports it; and \texttt{GestureType} being a proper bit-flag enum (its values are
powers of two, \texttt{Tap = 1} through \texttt{PinchComplete = 512}) means
\texttt{EnabledGestures} is genuinely a combinable set, not a single choice --- the
\texttt{|} composition in the worked example above is the ordinary, intended way to opt into
more than one gesture type at once.

\section{TextInputEXT and CNA-only NOXNA extensions}

\texttt{TextInputEXT} is not itself a CNA invention --- FNA ships this same extension beyond
the strict XNA~4.0 surface, and CNA ports it directly: \texttt{TextInput} / \texttt{TextEditing}
events, \texttt{StartTextInput} / \texttt{StopTextInput}, on-screen-keyboard visibility, and
input-rectangle placement. CNA layers three further, genuinely new members on top with no FNA
analog at all: IME composition-candidate lists (an SDL3-only capability), and an explicit
on-screen-keyboard/IME type hint.

Beyond input in the strict XNA sense, the \texttt{CNA::Input} namespace hosts several wholly
\texttt{NOXNA} subsystems already named in passing in this book's overview chapter:
\cnaclass{Clipboard} (system clipboard text); \cnaclass{Joysticks} (the raw, unmapped SDL
joystick API --- axes, buttons, hats, trackballs --- deliberately independent of
\cnaclass{GamePad}'s own mapped view of the same physical device); \cnaclass{Sensors} (host
device accelerometer/gyroscope, distinct from a gamepad's own motion sensors covered above);
\cnaclass{Power} (battery state); and \cnaclass{Haptics} (SDL3's own force-feedback API, with
a full effect-building lifecycle --- real actuation requires genuinely FFB-capable hardware,
since most ordinary gamepads only support simple rumble, not full force feedback).

\section{Platform-specific behavior worth knowing}

A handful of platform facts are worth stating plainly, because they are not CNA bugs even
though they look like gaps from inside a single platform: on Wayland,
\texttt{SDL\_GetGlobalMouseState} silently returns $(0,0)$, because the compositor's own
security model forbids querying the global pointer position at all --- absolute mouse
warping is similarly focus-gated, while relative (pointer-lock) mouse mode works normally.
On Windows, XInput controllers report no real USB vendor/product ID, so the \texttt{EXT}
GUID query returns the literal string \texttt{"xinput"} rather than a real hex identifier.
On Android and iOS, a touch device only reports as present \emph{after} its first touch ---
matching an identical quirk in FNA's own Windows/Android touch detection, not a CNA-specific
limitation. In a browser via Emscripten, gamepads remain invisible to SDL until the user
presses a button on one, a browser privacy restriction, and relative mouse mode maps onto the
Pointer Lock API, which itself requires a preceding user gesture to engage. Non-US keyboard
layouts introduce a genuine, structural XNA limitation rather than a CNA one: several
accented and non-ASCII characters common on European layouts have no corresponding
\cnaclass{Keys} enum value at all and are simply dropped in keycode mode --- \texttt{TextInputEXT}
is the correct API for capturing that text instead.

\section{Status}

At the time of writing, 524 input-specific tests pass in full, run shuffled five times over
to confirm order-independence. The fifteen formally blocked audit tasks are exclusively
real-hardware-only checks (physical keyboard/mouse/gamepad/touchscreen/IME verification, and
high-DPI display behavior) that cannot be automated at all, not gaps awaiting more
engineering effort. One acknowledged, low-priority documentation gap remains: the
\cnaclass{Joysticks} / \cnaclass{Sensors} / \cnaclass{Power} \texttt{NOXNA} extensions have solid
fake-backend unit coverage but no demo-UI or manual-hardware verification pass yet.
